% !TEX root = main.tex

\documentclass[../main.tex]{subfiles}
\begin{document}


\section{Subotamský mnich}
\label{sec:subatom}

\begin{povolani}[Subotamský mnich]
\textbf{Svolení}:\\
\textbf{Cena}: Zázraky nad rámec seslání zdarma se platí dovednostními body (viz Mechanismus zázraků).\\
\textbf{Popis}: Válečnické speciality a zázraky, viz níže.
\end{povolani}

\subsection{Válečnické speciality}
\label{sec:val-spec-sub}

\begin{specialita}[Bylinkářství]
\begin{itemize}
	\item \textbf{Vyžadovaný aspekt:} I
	\item \textbf{Popis:} Funguje jako standardní dovednost \dov{Medicína} - umožňuje sbírat byliny a z nich připravovat masti, výluh a další, zároveň lze použít jako improvizovanou první pomoc, je-li proto vhodné prostředí a narativ.
\end{itemize}
\end{specialita}

\begin{specialita}[Rychlé tasení]
\begin{itemize}
	\item \textbf{Vyžadovaný aspekt:} III
	\item \textbf{Prerekvizita:} Slušná Bojová dovednost, Dobrá \dov{Pozornost}, Dobrá \dov{Mobilita} nebo Slušná \dov{Zlodějina}
	\item \textbf{Popis:} Nikdo nepřátelský postavu nezastihne jinak, než se zbraní v ruce (pokud tedy nějakou zbraň má).
\end{itemize}
\end{specialita}

\begin{specialita}[Zteč]
\begin{itemize}
	\item \textbf{Vyžadovaný aspekt:} II
	\item \textbf{Prerekvizita:} Slušná \dov{Mobilita}
	\item \textbf{Popis:} Válečník se rychlým sprintem (nebo naopak šikovnými kličkami) rychle blíží ke střelci, a tím mu sníží počet výstřelů na jeho osobu o 1.
\end{itemize}
\end{specialita}

\begin{specialita}[Taneční parket]
\begin{itemize}
	\item \textbf{Vyžadovaný aspekt:} II
	\item \textbf{Prerekvizita:} \dov{mobilita} – dobrá
	\item \textbf{Popis:} Postava se může pokusit v nulovém kole přesunout do vedlejší zóny. Hází si standardně na past (akce \akc{Překonání}) a uspěje-li, v témže kole se do zóny přesune. Platí všechny výsledky akce \akc{Překonání}.
\end{itemize}
\end{specialita}

\begin{specialita}[Předvídavost]
\begin{itemize}
	\item \textbf{Vyžadovaný aspekt:} IV
	\item \textbf{Prerekvizita:} \dov{obezřetnost} – dobrá
	\item \textbf{Popis:} Šermíř nahlásí během boje dvě akce a při lámání chleba určí jednu z nich.
\end{itemize}
\end{specialita}

\begin{specialita}[Boj proti přesile]
\begin{itemize}
	\item \textbf{Vyžadovaný aspekt:} VI
	\item \textbf{Prerekvizita:} bojová dovednost – výborná, \dov{mrštnost} – výborná nebo \dov{odolnost} – výborná.
	\item \textbf{Další požadavky:} víc nepřátel
	\item \textbf{Popis:} Pro mistry boje přestávají platit obvyklé hranice. V případě, že bojuje proti více nepřátelům, nemohou použít Týmovou akci - musí s ním bojovat každý zvlášť.
\end{itemize}
\end{specialita}

\begin{specialita}[Sart]
\begin{itemize}
	\item \textbf{Popis:}
\end{itemize}
\end{specialita}


\subsection{Zázraky}
\label{sec:subotam-zazraky}

\subsubsection{Mechanismus zázraků}
Zázraky se používají za seslání. Kněz na každé úrovni aspektu obdrží ke každému zázraku nižší úrovně než je jeho dosavadní aspekt jedno seslání zdarma. Za každé další platí dovednostní. Zázraky se zpravidla sesílají v nulovém kole, není-li řečeno jinak.

\subsubsection{Krátkodobé modlitby (5)}

\begin{specialita}[Nápaditý boj]
\begin{itemize}
	\item \textbf{Seslání:} 0
	\item \textbf{Popis:} Subotam si obtiskne do mozku pozorované pohyby a získává tak možnost zkopírovat libovolnou fintu, kterou uvidí; tu pak smí sám použít. Může si jí taktéž držet v paměti do prvního použití; stejně tak ji může vypozorovat přímo v boji.
\end{itemize}
\end{specialita}

\begin{specialita}[Odhad soupeře]
\begin{itemize}
	\item \textbf{Seslání:} 0 pokud kněz vyhraje iniciativu/1
	\item \textbf{Popis:} Kněz zrakem svého boha prohlédne soka. Hodí si na svoji \dov{Pozornost} proti protivníkově \dov{Klamu}. V případě úspěchu pak odhalí jeho bojový aspekt a dokáže-li vhodným narativem popsat, jak ho využije pro získání výhody, má na něj jedno volné vyvolání. Lze použít na každého nepřítele jednou/souboj.
\end{itemize}
\end{specialita}

\subsubsection{Krátkodobé obřady (4)}

\begin{specialita}[Dorovnej šance]
\begin{itemize}
	\item \textbf{Seslání:} 0
	\item \textbf{Popis:} Bůh boje po prosbě kněze sejme negativní aspekt spojený s bojem z postavy (vlastní či cizí). Kněz si hodí na \dov{vůli} proti MÚ, která negativní aspekt na postavu umístila a vyhraje-li, aspekt je sejmut, je-li jednorázový.
\end{itemize}
\end{specialita}

\begin{specialita}[Neviditelnost]
\begin{itemize}
	\item \textbf{Seslání:} 1
	\item \textbf{Popis:} Kněz se zakryje roušku neviditelnosti. Na scénu se stane neviditelným; pro mechanické porovnávání má aspekt \asp{Neviditelný} a bonus +1 na relevantní dovednosti. V boji neviditelnost končí.

	Naopak, v průběhu boje může tímto zázrakem svoji zbraň učinit neviditelnou. Ta pak získává aspekt \asp{Smrt z neznáma} (bez volného vyvolání) a pasivní bonus +1 k hodnocení zbraně.
\end{itemize}
\end{specialita}

\subsubsection{Dlouhodobé modlitby (3)}

\begin{specialita}[Šampión bohů]
\begin{itemize}
	\item \textbf{Seslání:} 0
	\item \textbf{Popis:} Seš kněz boha války. Chápeš. Po seslání zázraku se kněžímu obnoví všechny políčka stresu a získá bonusové políčko drobného (+2) a mírného (+4) následku a k tomu se kněžímu neutralizují všechny DoT efekty jako \asp{Krvácení}, \asp{Otrava}. Trvá celou scénu.
\end{itemize}
\end{specialita}

\begin{specialita}[Heroismus]
\begin{itemize}
	\item \textbf{Seslání:} 0
	\item \textbf{Popis:} Pořád seš kněz boha války. V případě, kdy osudí není na tvé straně - nepřátel je mnohem více, jste v bezradné situaci apod. (podrobnosti na uvážení PJ), Gor si uslintne a dá ti, co potřebuješ. Získáváš aspekt \asp{Heroismus} s jedním volným vyvoláním, který je použitelný např. pro strhnutí davu, zastrašení apod., bonus +1, který smíš do konce scény přiřadit libovolné dovednosti, a nakloníš si osud zpět k sobě - získáváš 2 body osudu.
	\item \textbf{V případě krajní nouze:} Subotam se rozhodne, že toto je vhodná chvíle skončit. Sešle Heroismus, a dokud svůj úkol nesplní, nemůže zemřít. Je na rozhodnutí PJ, zda-li se souboj odehraje mechanicky, či jestli bude zajímavější jej pouze odvyprávět. V případě, že se Subotam pro tuto možnost rozhodne, nesmí používat žádné Gorovy zázraky (jeho přízeň si vyčerpal), dokud jej osobně nenavštíví v jeho zahradě ve Stínovém světě.
\end{itemize}
\end{specialita}

\subsubsection{Dlouhodobé obřady (2)}

\begin{specialita}[Požehnání v boji se zlem]
\begin{itemize}
	\item \textbf{Seslání:} 1 (mimo boj)
	\item \textbf{Popis:} Subotam provede rituál a připraví se a své spolubojovníky tak na nadcházející souboj s temnou stranou síly. Každý zúčastněný dostane bonus +2, který smí přerozdělit do posílení svého Hodnocení zbraně či zbroje. Trvá celou scénu.
\end{itemize}
\end{specialita}

\begin{specialita}[Posílení vlastnosti]
\begin{itemize}
	\item \textbf{Seslání:} 1 (mimo boj)
	\item \textbf{Popis:} Subotam provede rituál a připraví se a své spolubojovníky tak na nadcházející souboj s temnou stranou síly. Každý zúčastněný dostane bonus +2, který smí přerozdělit do posílení svých dovedností.
\end{itemize}
\end{specialita}


\end{document}



%%% Local Variables:
%%% mode: LaTeX
%%% TeX-master: "../../main"
%%% End:
